\section{Procedure}
\label{sec:study_design_procedure}

Each session will begin with participant information, consent, and a short health and comfort check. The participant will then be introduced to the rowing setup and the VR headset in neutral terms. The experimenter will explain that the study compares two virtual representations of the same rowing action and that the participant should report their own experience rather than trying to identify a correct or preferred technical solution.

Before the measured trials, the system will be calibrated to the participant and the physical rowing setup. The participant will then complete both representation conditions. In each condition, the participant will first receive a short familiarization period, followed by a standardized rowing period. The familiarization period is included because movement-based ownership can take time to emerge, and very short exposures may capture initial confusion rather than a stabilized experience of the mapping \parencite{kalckertEhrsson2017ownershipOnset}.
% SOURCE-CHECK: MD Papers/kalckert_ehrsson_2017_moving_rubber_hand_ownership_onset_time_markdown_bundle/kalckert_ehrsson_2017_moving_rubber_hand_ownership_onset_time_thesis_notes.md

\drafttodo{TODO:CONFIRM exact familiarization duration and exact rowing duration per condition. Current default: at least one minute of stable familiarization followed by a brief standardized rowing period.}

After each condition, the participant will complete the per-condition questionnaire. The questionnaire will include ownership and agency items, exploratory body-location items, and thesis-specific items on embodied action congruence, the real handle as a physical anchor, virtual hand/oar mapping, sport-action plausibility, and perceived mapping mismatch. Safety and comfort symptoms will be recorded separately from these construct ratings.

A short rest will be offered between conditions. After both conditions have been completed, the participant will answer comparative preference questions and may provide open comments about the main differences they noticed. The participant may stop the session at any time, and the experimenter will stop the session if the participant reports nausea, dizziness, pain, distress, or unsafe discomfort.
